% Vietnamese translation for OI Public Library
% Source-File: IOI中国国家候选队论文/2006/龙凡/龙凡.ppt
\documentclass[11pt]{article}
\usepackage{oipl-vn}

\title{Một lớp bài toán đoán số\\{\large Ghi chú theo slide}}
\author{Long Fan}
\date{2006}

\begin{document}
\maketitle

\origtitle{Number-guessing problems}
\sourcepath{IOI China candidate team papers / 2006 / Long Fan / Long Fan.ppt}

\section{Phạm vi}

Bộ trình chiếu đi kèm bài viết về các biến thể của bài đoán số. Ghi chú này
dựa trên bản bài viết đã dịch và giữ các truy hồi chính.

\section{Bài cơ bản}

Nếu mỗi lần hỏi nhận ngay câu trả lời nhỏ hơn, bằng hoặc lớn hơn, chiến lược
quen thuộc là tìm kiếm nhị phân. Nhìn theo khả năng bao phủ, gọi \(f(i)\) là
độ dài đoạn lớn nhất xử lý được bằng \(i\) câu hỏi. Khi đó
\[
  f(0)=0,\qquad f(i)=2f(i-1)+1.
\]

\section{Giới hạn một hướng trả lời}

Nếu số lần nhận câu trả lời \(X>Y\) bị giới hạn, trạng thái cần thêm số lần
còn được phép nhận hướng đó. Gọi \(f(i,j)\) là độ dài lớn nhất xử lý được bằng
\(i\) câu hỏi khi còn \(j\) lần cho hướng \(X>Y\). Khi hỏi một giá trị, phía
nhỏ hơn không tiêu thụ hạn mức, phía lớn hơn tiêu thụ một lần:
\[
  f(i,j)=f(i-1,j)+1+f(i-1,j-1).
\]

\section{Độ trễ phản hồi}

Khi câu trả lời của câu hỏi hiện tại đến muộn một bước, chiến lược phải hỏi
trong lúc chưa biết kết quả trước đó. Bài viết xử lý bằng cách mở rộng trạng
thái để chứa thông tin đang chờ và tính khả năng bao phủ trong trường hợp xấu
nhất.

\section{Ghi nhớ}

Thay vì cố sửa trực giác nhị phân cho từng biến thể, hãy mô tả năng lực của
chiến lược bằng truy hồi. Khi các ràng buộc thay đổi, trạng thái truy hồi thay
đổi theo nhưng cách chứng minh vẫn giữ cùng tinh thần.

\end{document}
